\begin{table}[t]
\tbl{\label{tab:povray2}}
{\begin{tcolorbox}[tabularx={p{3cm}|p{2cm}|p{6.3cm}|},enhanced,fonttitle=\bfseries\large,fontupper=\normalsize\sffamily,
colback=yellow!10!white,colframe=red!50!black,colbacktitle=blue!30!white,
coltitle=black,title=Povray Interface (Part 2)]
\hline
Command & Arguments & Purpose \\
\hline
\hline
\verb'-clipping'  &  $r$ $c$  &  In round $r$, set clipping radius to $c$. \\
\hline
\verb'-text'  &  $r$ $a$ text  &  In round $r$, produce running text \verb'text'  with sustain value $a$. \\
\hline
\verb'-label'  &  $r$ $s$ $a$ $g$ text  &  In round $r$, produce running text \verb'text' with start value $s$, sustain $s$ and gravity $g$.  \\
\hline
\verb'-latex'  &  $r$ $s$ $a$ praeamble $g$ text $l$ fname  &  In round $r$, produce running latex text \verb'text' with start value $s$, sustain $s$ and gravity $g$. Put \verb'praeamble' in the latex source code. Use \verb'fname' for the latex file names (no extension). \\
\hline
\verb'-global_picture_scale'  &  $d$  &  Set scaling factor to $d$. \\
\hline
\verb'-picture'  &  $r$ $d$ fname options  &  In round $r$, place picture \verb'fname' scaled by $d$ using options. \\
\hline
\verb'-picture'  &  $r$ $d$ fname options  &  In round $r$, place picture \verb'fname' scaled by $d$ using options. \\
\hline
\verb'-look_at'  &  $L$  &  Override camera look-at value to $L$. $L$ is a three-dimensional vector. \\
\hline
\verb'-default_angle'  &  $a$  &  Set default camera angle to $a$. \\
\hline
\verb'-clipping_radius'  &  $f$  &  Set default clipping radius to $f$. \\
\hline
\verb'-scale_factor'  &  $s$  &  Set default scale factor to $s$. \\
\hline
\verb'-line_radius'  &  $s$  &  Set default line radius to $s$. \\
\hline
\end{tcolorbox}}
\end{table}
